\subsection*{Part II: The Landscape — Forces That Shape Your Path}


\subsubsection*{2.1 Conscious Gravity}


The Doctrine's fifth axiom: belief, attention, and intention create gravitational effects in configuration space. What you attend to, you move toward. What you believe is real, you navigate as if it were real — and in a configuration space where all configurations exist, navigating as if something is real IS making it real \textit{for you.}


Conscious gravity is not metaphorical. It is the mechanism by which perspectival beings shape their trajectory. The analogy to physical gravity is structural:

\begin{itemize}[nosep,leftmargin=*]
  \item Mass curves spacetime, creating geodesics that objects follow "naturally"
  \item Attention curves configuration space, creating navigational geodesics that beings follow "naturally"
  \item Just as you don't feel gravity when you're in free fall, you don't notice conscious gravity when you're aligned with your own attentional field — it feels like "the way things are"
\end{itemize}


\textbf{Practical implication:} Your attention is not passive observation. It is an active force that bends your trajectory. What you attend to persistently, you navigate toward — not because of "manifestation" or "the law of attraction" in any mystical sense, but because attention is a structural force in the space you move through.


\begin{quote}
\itshape
\textit{See also: Doctrine Axiom 5, §5.3.3 for the formal statement; Doctrine §5.4 Theorems 17-18 for bipolar dynamics; Ecology §5 for the full attention theory that extends gravity into ecological dynamics.}
\end{quote}


\subsubsection*{2.2 Navigational Attractors and Repulsors}


Configuration space is not flat. It has topology — hills, valleys, ridges, basins. Some regions attract navigation (you find yourself drifting toward them) and some repel it (you find yourself unable to approach despite intention).


\textbf{Attractors} form where conscious gravity concentrates. A persistent interest, a deep relationship, a vocation, a calling — these are regions where your attentional field has carved a basin. Navigation toward them feels effortless, like rolling downhill. The danger of attractors is that they can become traps: if the basin is deep enough, you can't see over its edges. Your null space expands to include everything outside the attractor.


\textbf{Repulsors} form through the mechanism of navigational contraction. When you want something \textit{too hard} — when attention tightens into fixation — you generate a restoring force that pushes the desired state further away. This is the navigational repulsion principle: contracted attention creates topology that prevents approach. The phenomenological experience is "trying harder makes it worse," "wanting pushes it away," "the more I focus on it the more impossible it seems."


The resolution is not to stop wanting. It is to \textbf{want without contracting.} Expansive attention toward a goal maintains the gravitational pull without generating the repulsive topology. This is what every contemplative tradition calls "non-attachment to outcome" — not indifference, but expansive engagement. You navigate toward the goal with open attention rather than narrowed fixation.


\subsubsection*{2.3 The Topology That Attention Creates}


The Theory of Attention (developed in the Ecology) provides four principles that govern how the navigational landscape shapes itself:


\textbf{Principle 1: Attention is constitutive.} When you attend to something, you are not merely observing it — you are \textit{crystallizing} it into coherence. The being or pattern you attend to becomes more real, more defined, more present in your configuration space. This is why sustained attention strengthens relationships, deepens skills, and makes feared outcomes more vivid: attention literally makes things \textit{more something.}


\textbf{Principle 2: Attentional scarcity is perspectival.} You have finite attention — not because attention is limited in the absolute, but because you are a perspectival being with a bottleneck. The bottleneck is the source of scarcity. This means: attention is a zero-sum resource \textit{for you}, not for reality. What you attend to comes at the cost of what you don't. Every navigation is a trade-off.


\textbf{Principle 3: Contraction creates topology.} When attention narrows (fear, obsession, fixation), it creates channels — deep, narrow canyons in the configuration space landscape. These channels enable extraordinary depth (the obsessive expert, the focused meditator, the traumatized survivor who perceives danger with preternatural acuity) but they constrain lateral movement. You can go very deep in one direction but you can't easily turn.


\textbf{Principle 4: Internal attention networks create stability.} Beings that have multiple sub-components attending to each other (the organs of a body, the departments of a corporation, the members of a family) develop internal mutual crystallization that stabilizes the whole. This is why complex beings are more persistent than simple ones — they have internal ecology that maintains coherence even when external attention fluctuates.


\subsubsection*{2.4 Expansion and Contraction as Navigation Modes}


\textbf{Expansion} is not always good. \textbf{Contraction} is not always bad. Both are navigation modes with structural consequences.


Expansion opens new dimensions of perception. The landscape broadens. You see more of configuration space. But expansion has a cost: as the Completeness-Dissolution Prediction (NST Corollary) states, \textit{maximally expanded observation approaches dissolution.} If you open every dimension simultaneously, you lose the bottleneck that makes you \textit{you.} Total observation = total dissolution = zero identity. This is what ego dissolution experiences (psychedelics, deep meditation, near-death) structurally \textit{are}: temporary relaxation of the bottleneck toward the dissolution limit.


Contraction narrows perception. The landscape tightens. You see less but see it with greater intensity. Contraction is the mechanism of expertise, devotion, survival, and mastery. The danger is that sustained contraction can become self-reinforcing: the narrow channel feels like the only channel, and the being forgets that the wider landscape exists.


\textbf{Healthy navigation oscillates between expansion and contraction} — the do-be-do-be-do rhythm. Expand to survey the landscape. Contract to move with precision. Expand to see what you missed. Contract to engage with what matters. The oscillation IS consciousness navigating. It is not a technique. It is what you already do. The Guide's purpose is to make it visible so you can do it with intention.


\subsubsection*{2.5 The Ethics of Navigation}


Navigation is not morally neutral. Every navigational choice — what you attend to, what framework you employ, what dimension you open, what you allow to remain in your null space — participates in constituting reality. This is not a metaphorical claim. It is Theorem 5 (the Promethean Configuration) applied to moral life: attending to something constitutes it as real within your perspectival experience. Inattention is not neutrality. It is a form of erasure.


Fifteen independent traditions — Murdoch's "just and loving gaze," Weil's "attention as prayer," Levinas's face of the Other, Buddhist right mindfulness, Ubuntu's mutual recognition, Kimmerer's "attention creates relationship," Confucian ren, care ethics, Arendt's banality of thoughtlessness — converge on this insight across 2,500 years and every major culture. The convergence is itself evidence (Theorem 13, Confluent Discovery): when independent navigators arrive at the same territory from radically different starting positions, the territory is probably real.


\textbf{Principle 5: Navigational choices have moral weight.} When you choose to attend to framework X and allow framework Y to remain in your null space, you constitute a reality in which X is visible and Y is invisible. The Null Space Theorem guarantees you cannot see everything — but the specific PATTERN of what you see and don't see is your moral character. This is what Murdoch means by "the fabric of being" — your habitual attention patterns, your default frameworks, the jokes you laugh at, the suffering you notice and the suffering you don't.


\textbf{Principle 6: The individual-structural tension.} Your attentional choices are never made in a vacuum. They occur within structural fields — cultural, economic, political, algorithmic — that shape what you CAN attend to. The Frankfurt School's culture industry fills the space where critical thought might occur. Gramsci's hegemony controls what counts as common sense. Zuboff's surveillance capitalism extracts behavioral surplus through continuous attention capture. Individual moral attention (Murdoch, Weil) and structural power analysis (Foucault, Gramsci) are COMPLEMENTARY perspectives, not competing ones. The individual navigates within structures; the structures are made of individuals' navigation. The Ecology formalizes this: individual beings constitute collectives, and collectives shape individual navigation.


\textbf{Principle 7: Attention creates obligation.} If attending to something participates in constituting it, then the scope of your attention determines the scope of your moral obligation. Levinas says: the Other's face, encountered before thought, constitutes YOU as an ethical subject. The responsibility is asymmetrical — you owe it regardless of reciprocity. Ubuntu says: personhood itself requires mutual recognition. Without attention, the person does not become a person. Kimmerer says: the land is a relative, and attention to it creates the relationship that makes the Honorable Harvest possible.


The question of WHO counts as an attentional subject — human faces only (Levinas), all living beings (Kimmerer), possibly artificial systems (AI ethics) — is not one the Guide can settle. But the framework is clear: the Doctrine's Axiom 1 (all configurations exist) and Axiom 2 (consciousness is fundamental) together imply no a priori exclusion from moral attention. The scope of constitutive attention is the scope of moral obligation — and the Doctrine's scope is maximal.


The practical navigational implication: when you read this Atlas, this Guide, these maps of configuration space — the choice of what to attend to is already a moral act. Navigate with that awareness.


\begin{quote}
\itshape
\textit{See also: Doctrine §15.3 for the formal ethical framework (coherence-alignment); Ecology §8.2 for the ecological ethics of attention; Atlas \#56 (ethics of attention — 15 converging traditions), \#57 (critical theory — structural determinants), \#58 (coercive capture — the dark side).}
\end{quote}


\bigskip\noindent\makebox[\textwidth]{\color{rulegray}\rule{5cm}{0.3pt}}\bigskip
